
\section{Detektorempfänger}
\label{section:detektorempfänger}
\begin{frame}%STARTCONTENT
    \begin{columns}
        \begin{column}{0.48\textwidth}
            \begin{itemize}
                \item Einer der ersten und einfachsten Empfänger für AM
                \item Energie wird direkt aus dem empfangenen Signal gezogen
                \item Nur im Lokalbereich von starken Rundfunksendern nutzbar
            \end{itemize}
        \end{column}
        \begin{column}{0.48\textwidth}
            \begin{figure}
                \includegraphics[width=0.8\textwidth]{foto/193}
            \end{figure}
        \end{column}
    \end{columns}
\end{frame}

\begin{frame}{Der einfachste AM-Empfänger}
    \begin{columns}[T]
        \begin{column}{0.48\linewidth}
            \DARCimage{0.85\linewidth}{799include}

            \vspace{1cm}
            \href{https://tinyurl.com/2yrgdghz}{\textcolor{DARCblue}{\faLink~Simulation}}
        \end{column}
        \begin{column}{0.48\linewidth}
            \visible<2->{\DARCimage{0.85\linewidth}{800include}}

            \vspace{10pt}

            \visible<3->{\DARCimage{0.85\linewidth}{801include}}

            \visible<4->{\DARCimage{0.85\linewidth}{802include}}
        \end{column}
    \end{columns}
\end{frame}


\begin{frame}
    \only<1>{
        \begin{PQuestion}{EF101}{Was stellt nachfolgende Schaltung dar?}{Modulator}
            {Verstärker}
            {Oszillator}
            {Detektorempfänger}
            {\DARCimage{1.0\linewidth}{523include}}\end{PQuestion}

    }
    \only<2>{
        \begin{PQuestion}{EF101}{Was stellt nachfolgende Schaltung dar?}{Modulator}
            {Verstärker}
            {Oszillator}
            {\textbf{\textcolor{DARCgreen}{Detektorempfänger}}}
            {\DARCimage{1.0\linewidth}{523include}}\end{PQuestion}

    }
\end{frame}%ENDCONTENT

\begin{frame}{Geradeausempfänger}
    \begin{Columns}[T]
        \Column{0.5}{
            \DARCimage{\linewidth}{736include}
            \vspace{1cm}
        }
        \Column{0.5}{
            \begin{itemize}[<+->]
                \item Das von der Antenne aufgenommene Signal bleibt in seiner Frequenz bis zur Demodulation erhalten.
                \item Ein einzelner Filter hat zu geringe Trennschärfe
                \item Deshalb werden mehrere Filter kombiniert um eine höhere Trennschärfe zu erreichen
                \item Alle Filter müssten dann exakt auf die Empfangsfrequenz angepasst werden, dies ist Aufwändig
            \end{itemize}
        }
    \end{Columns}
\end{frame}